% ============================================================
%  Chapter 4: Design and Implementation
% ============================================================
%
% AIM: Describe HOW you built the system. A reader should be able to
%      reproduce your work from this chapter alone. ~6-8 pages.
%
% CONTENTS:
%   4.1 Software Framework (Why LION, Why tomosipo, operator abstraction)
%   4.2 Dataset Pipeline (which dataset, preprocessing, splits)
%   4.3 Model Architecture (denoiser U-Net, policy network, value network)
%   4.4 Training Procedure (denoiser pre-training, RL loop, losses, hardware)
%   4.5 Baseline Implementations (FBP, TV, SIRT, LPD)
%
% FIGURES:
%   - Fig 4.1: System architecture diagram [TikZ]
%   - Fig 4.2: U-Net denoiser architecture
%   - Fig 4.3: Policy network architecture
%   - Fig 4.4: Training curve (reward vs episode) [generate during training]
%
% TIPS:
%   - Include code snippets where they clarify
%   - Be specific about dimensions, channels, hyperparameters
%   - Write this AS you implement, not after
%   - Mention any deviations from the original paper and why
% ============================================================

\chapter{Design and Implementation}
\label{chap:implementation}

% -----------------------------------------------------------
\section{Software Framework}
\label{sec:software-framework}
% -----------------------------------------------------------

\subsection{Why LION?}

LION (\emph{Learned Iterative Optimisation Networks}) is an open-source Python toolbox
developed by the Cambridge Computational Imaging (CIA) group at DAMTP, available at
\url{https://github.com/CambridgeCIA/LION} \citep{Biguri2016TIGRE}. LION was chosen as
the implementation framework for this project for several reasons:

\begin{itemize}
    \item \textbf{Integrated CT pipeline.} LION provides end-to-end infrastructure for
          CT reconstruction research: dataset loading (including the Mayo Clinic Low Dose
          CT Grand Challenge dataset \citep{MayoClinic}), geometry configuration, sinogram
          simulation, classical reconstruction baselines (FBP/FDK, SIRT, TV), and learned
          reconstruction model training loops.
    \item \textbf{PyTorch native.} LION is built on PyTorch, allowing seamless integration
          of differentiable components---a prerequisite for the model-based RL training in
          TFPnP \citep{Wei2022TFPnP}, which requires backpropagation through the PnP
          iterations.
    \item \textbf{Modular solver architecture.} LION's \texttt{Solver} abstraction
          separates the optimisation loop from the model definition, making it
          straightforward to implement a custom TFPnP training loop alongside the existing
          \texttt{SupervisedSolver} used by LPD \citep{AdlerOktem2018}.
    \item \textbf{TIGRE/tomosipo integration.} LION uses tomosipo and TIGRE under the
          hood, providing cone-beam CT operators not available in TorchRadon (which the
          original TFPnP paper used).
\end{itemize}

\subsection{Why Tomosipo?}

Tomosipo \citep{Hendriksen2021tomosipo} is a Python library for defining CT acquisition
geometries and computing forward/back-projection operators via GPU-accelerated ASTRA
Toolbox backends. Its role in this project is:

\begin{itemize}
    \item \textbf{Geometry abstraction.} A clean API for specifying parallel-beam and
          cone-beam scan geometries:
          \begin{verbatim}
import tomosipo as ts
vg = ts.volume(shape=(1, 256, 256))
pg = ts.parallel(angles=30, shape=(1, 256))
A  = ts.operator(vg, pg)
          \end{verbatim}

    \item \textbf{Operator abstraction.} The resulting \texttt{A} object implements
          $\bm{x} \mapsto A\bm{x}$ (forward projection) and $\bm{y} \mapsto A^T\bm{y}$
          (backprojection) as callable functions.

    \item \textbf{Differentiability.} Via \texttt{ts\_algorithms}, the projectors can be
          wrapped in PyTorch's autograd framework.

    \item \textbf{Consistency with the paper's geometry.} The TFPnP paper uses a 30-view
          parallel-beam geometry with 182 detector pixels. Tomosipo can replicate this
          exactly.
\end{itemize}

\subsection{Operator Abstraction}

At the core of the CT pipeline is the system operator pair $(A, A^T)$. In LION/tomosipo,
these are implemented via the ASTRA Toolbox's CUDA-accelerated routines. The forward
projection $\bm{y} = A\bm{x}$ simulates the physical measurement process described in
\cref{sec:ct-basics}. The backprojection $A^T\bm{y}$ distributes each sinogram value
back along its ray into the image volume.

\subsection{Forward and Backprojection Implementation}

For the $\bm{z}$-step of ADMM with quadratic data fidelity
$D(\bm{y}, A\bm{x}) = \frac{1}{2}\|A\bm{x} - \bm{y}\|_2^2$, the proximal operator
has the closed-form solution:
\begin{equation}
    \operatorname{Prox}_{\frac{1}{\mu}D}(\bm{v})
        = \left(\mu A^T A + I\right)^{-1}\!\left(\mu A^T \bm{y} + \bm{v}\right),
    \label{eq:data-prox}
\end{equation}
which requires an iterative solver. In the TFPnP implementation \citep{Wei2022TFPnP},
a single gradient step is used as a tractable approximation:
\begin{equation}
    \bm{z}^{k+1} \approx \bm{v} - \frac{1}{\mu} A^T\!\left(A\bm{v} - \bm{y}\right),
    \label{eq:data-grad-step}
\end{equation}
where $\bm{v} = \bm{x}^{k+1} + \bm{u}^k$. This is implemented as two tomosipo operator
calls and is fully differentiable.

\subsection{Tomosipo Exploration and Validation}
\label{sec:tomosipo-exploration}

% AIM: Justify WHY you did a hands-on tomosipo exploration before
% implementing TFPnP, and explain what each experiment taught you.
% This maps to Ander's Milestone 1: "Can you define a geometry,
% forward project an arbitrary image, add noise, and reconstruct?"

Before implementing TFPnP, a systematic exploration of tomosipo was carried out
(documented in \texttt{notebooks/01\_tomosipo\_fundamentals.ipynb}) to build
operational understanding of the CT simulation pipeline. This step was essential for
two reasons: first, the TFPnP paper uses TorchRadon \citep{Ronchetti2020} as its CT
backend, whereas this project uses tomosipo/ASTRA; verifying that the same physical
operations (forward projection, backprojection, noise simulation) behave equivalently
in the new backend is a prerequisite for a valid reproduction. Second, every component
of the TFPnP algorithm---the data-consistency $\bm{z}$-step~\eqref{eq:data-grad-step},
the reward function, and the training loop---depends on correct and efficient use of
$A$ and $A^T$, so fluency with the operator interface is necessary before writing any
model code.

The exploration proceeded in three stages:

\paragraph{Stage 1: Geometry and operator fundamentals.}
Cone-beam and parallel-beam geometries were defined and compared. Cone-beam geometry
(\texttt{ts.cone}) was used initially for general 3D exploration, as it represents the
most common clinical CT configuration and exercises the full tomosipo API. The
exploration then moved to 2D parallel-beam geometry (\texttt{ts.parallel} with a
single detector row), which matches the TFPnP experimental setup: 30 projection angles
over an $N \times N$ image with $N$ detector pixels. The distinction matters because
parallel-beam projections are shift-invariant (enabling analytical FBP inversion via
the Central Slice Theorem), whereas cone-beam projections are not, requiring the FDK
approximation \citep{Feldkamp1984}. All subsequent project experiments use parallel-beam
geometry to maintain consistency with the reference paper.

\paragraph{Stage 2: SIRT implementation from scratch.}
The SIRT algorithm was implemented manually using the tomosipo operator interface,
following the update rule from \cref{sec:classical-reconstruction}. This served two
purposes: (i)~it verified that the forward operator $A$ and its adjoint $A^T$ behave
correctly (the reconstruction converges to the ground truth for noiseless, fully sampled
data); and (ii)~it confirmed that the GPU-accelerated PyTorch path through tomosipo
produces identical results to the NumPy path, which is important since TFPnP training
requires GPU tensors throughout. The manual implementation was validated against
\texttt{ts\_algorithms.sirt} and ASTRA's internal \texttt{SIRT3D\_CUDA}, with negligible
differences in all cases.

\paragraph{Stage 3: Noise simulation and baseline reconstruction.}
Gaussian noise at the three levels used in the TFPnP paper ($\sigma \in \{5\%, 7.5\%,
10\%\}$) was added to sparse-view (30-angle) sinograms. Four reconstruction algorithms
were then compared: FBP (via \texttt{tsa.fbp}), SIRT, Nesterov-accelerated least squares
(NAG\_LS), and TV minimisation (via \texttt{tsa.tv\_min2d}). PSNR and SSIM were computed
for each, establishing the baseline performance landscape against which TFPnP will be
evaluated. The key observations were: (i)~FBP from 30 views produces severe streak
artefacts and low PSNR, confirming the ill-posedness described in
\cref{sec:ill-posedness}; (ii)~TV minimisation substantially outperforms unregularised
methods but introduces staircase artefacts; (iii)~the quantitative gap between TV and
the TFPnP results reported in \cref{tab:ct-results} defines the improvement that a
learned method must achieve. These baseline reconstructions also serve as the reference
images for the qualitative comparison figures in \cref{chap:results}.


% -----------------------------------------------------------
\section{Dataset Pipeline}
\label{sec:dataset-pipeline}
% -----------------------------------------------------------

% TODO: Fill in once you have decided on dataset and implemented loading.


% -----------------------------------------------------------
\section{Model Architecture}
\label{sec:model-architecture}
% -----------------------------------------------------------

% TODO: Fill in once you have implemented the model.


% -----------------------------------------------------------
\section{Training Procedure}
\label{sec:training-procedure}
% -----------------------------------------------------------

% TODO: Fill in during/after training.


% -----------------------------------------------------------
\section{Baseline Implementations}
\label{sec:baselines}
% -----------------------------------------------------------

% TODO: Describe how you ran the baselines.